\section*{Acknowledgements}
This paper makes use of the following ALMA data: the \NProjCodes{} projects
listed in Table~\ref{tab:projcodes}, each cited as \texttt{ADS/JAO.ALMA\#}
followed by its project code. The \NEB{} primary-census execution blocks
searched span \PrEbDateFirst{} to \PrEbDateLast{}, and we thank the principal
investigators and teams who proposed and obtained
them.\ProjCodeUnresClause{} ALMA is a partnership of ESO
(representing its member states), NSF (USA) and NINS (Japan), together with NRC
(Canada), NSTC and ASIAA (Taiwan), and KASI (Republic of Korea), in cooperation
with the Republic of Chile. The Joint ALMA Observatory is operated by ESO,
AUI/NRAO and NAOJ. This work has made use of data from the European Space
Agency (ESA) mission {\it Gaia} \citep{GaiaDR3}, processed by the {\it Gaia}
Data Processing and Analysis Consortium (DPAC). This research has made use
of the SIMBAD database \citep{SIMBAD2000} (CDS, Strasbourg, France); of the
NASA Exoplanet Archive \citep{NEA2013}, which is operated by the California
Institute of Technology under contract with the National Aeronautics and
Space Administration under the Exoplanet Exploration Program; and of the
Splatalogue spectral-line database \citep{Splatalogue2007}.

\emph{Use of generative artificial intelligence.} Large language models were
used in this work. They wrote most of the software: the archival retrieval and
calibration scripts, the search and vetting pipeline, the statistical tests,
and the generators that produce every number, table and figure in this paper.
They also drafted text, which the authors revised. The authors set the
scientific question, chose every criterion, and are responsible for the
content of this paper.

\emph{Testing and validation of that analysis.} No result here is accepted
on a model's assertion. No number in this paper is typed: each is written by
the program that computes it from the data, so a value in the text cannot
differ from the analysis behind it. A separate audit re-derives the
identities that must hold among the published values -- counts that must
partition, quantities that must sum, numbers reached by two routes that must
agree -- and refuses the paper where one does not. Every assertion in the
analysis was first made to fail on purpose, on deliberately wrong input,
before it was relied on, since an assertion that has never fired is no
evidence that anything is right. Every reference cited here has been checked
against its source of record.

\emph{Software.} The pipeline depends on \textsc{casa}~6.7.6
\citep{CASA2022}, \textsc{python}~3 with \textsc{astropy} \citep{Astropy2022}, \textsc{numpy}
\citep{Harris2020} and \textsc{scipy} \citep{Virtanen2020}, and the ALMA Science
Archive and SIMBAD TAP interfaces through \textsc{astroquery} \citep{Ginsburg2019};
all are free and open-source, at the versions cited.

\emph{Author contributions.} G.J.W.\ conceived the survey, defined the target sample and the selection
criteria, and wrote the manuscript. R.D.\ built and operated the archival
retrieval, calibration and search pipeline, implemented the spatial control
statistic and the injection--recovery campaigns, and produced the per-window
catalogue and its regeneration scripts. Both authors designed the vetting
procedure, chose the masked species, carried out the statistical
validation reported in \S\ref{sec:statistic} and
Appendix~\ref{app:falsealarm}, and agreed the dispositions of
Table~\ref{tab:ledger}.
%% Authors to confirm this division before submission.

\emph{Funding.} This work received no dedicated funding. It used public archival data and
no observing time was awarded for it.
%% Authors to replace this if any grant supported the work.

\emph{Competing interests.} The authors declare no competing interests. VBRL Holdings is a privately
held company through which R.D.\ conducts independent research, named
here solely as his affiliation; it had no role in the design of the
survey, the selection of targets, the analysis, the decision to publish
or the content of this paper, and provided no funding for the work.
%% Authors to confirm this wording, which is accurate only if the company's
%% involvement was indeed nil.

\section*{Data availability}
\label{sec:data}
Every observation this paper rests on is public, and is identified here
precisely enough to be fetched. The \NEB{} primary-census execution blocks and
the blocks of every other set named in Table~\ref{tab:blockledger} are held in
the ALMA Science Archive under the \NProjCodes{} project codes of
Table~\ref{tab:projcodes}, and are retrieved there by the execution-block
identifier, which Table~\ref{tab:ledger} gives for every one of the
\EvNCrossRaw{} threshold crossings. ALMA's own citation
convention for archival data is the project code rather than a persistent
identifier, so the archive's identifiers are given in place of one, which is
the recommended practice where a data set carries none of its own
\citep{Chen2022}. The
ancillary catalogues are public in the same sense: Gaia~DR3 for the astrometry
and the colours, the NASA Exoplanet Archive for the planet counts, and the
CDMS and JPL line catalogues through Splatalogue for the mask, queried once on
\McQueryDate.

The derived results are in the paper. Table~\ref{tab:blockledger} partitions
the blocks, windows, stars, systems and hours of every set the search defines;
Table~\ref{tab:counts} carries every headline count, the five fates of the
unprocessed blocks and the crossing partition; Appendix~\ref{app:persys} and
Table~\ref{tab:persys} give the sensitivity, the frequency, the channel
width, the covered bandwidth and the confirmability of every system of the
carrier search;
Table~\ref{tab:ledger} lists all \EvNCrossRaw{}
crossings with their frequencies in both frames, the stellar and control
statistics, the nearest masked transition and the recurrence test;
Table~\ref{tab:cand} gives the two leading events in full; and
Table~\ref{tab:holdout} does the same for the reserved hold-out.

%% ★★★ THE MACHINE-READABLE FORM OF TWO PUBLISHED TABLES (referee 2, minor 11),
%% as arXiv ANCILLARY FILES -- `anc/` in the submission package.  This journal
%% publishes the arXiv PDF itself and hosts no supporting information, so an
%% ancillary file is the mechanism available, and it is the right one: arXiv's
%% own first listed use for it is "raw data for tables and plots in the
%% article" (info.arxiv.org/help/ancillary_files.html).
%% ★★ WHAT THIS PARAGRAPH MAY PROMISE IS LIMITED BY WHAT THE FILES HOLD, and
%% that is why the two column lists are MACROS: `si_r18.py` builds
%% \SiCatColList and \SiLdgColList phrase by phrase from the headers of the
%% files on disk and asserts that the phrases cover those headers exactly and
%% in order, so this sentence cannot promise a column a file does not carry.
%% The row counts are the paper's own published counts, asserted against the
%% files in the same generator (S1, S7).
%% ★ An ancillary file is not refereed and is not part of the article.  So
%% nothing here is offered as evidence, no claim rests on it and no identifier
%% of ours is quoted.  ★★ AND THE TWO FILES ARE NOT THE SAME KIND OF THING,
%% which the sentence must not blur: the ledger is the machine-readable form of
%% a table this paper typesets in full, while the per-window catalogue is the
%% window-level measurement the paper's counts are formed from and is typeset
%% nowhere -- \NWindows{} rows would be forty pages.  Table~\ref{tab:persys} is
%% per SYSTEM, from each system's best window, so it is not that file either.
%% A first draft of this paragraph said both files held published tables; it
%% was wrong about the catalogue, and the distinction is the whole value of
%% saying only what the files contain.
Two of these tables accompany the paper on the arXiv in machine-readable form,
as ancillary files. The per-window catalogue carries one row for each of the
\NWindows{} searched spectral windows, with \SiCatColList. The crossing ledger
carries one row for each of the \EvNCrossRaw{} threshold crossings, with
\SiLdgColList. Each is plain tab-separated text and describes its own columns
in its own header. The ledger file holds the rows of
Table~\ref{tab:ledger} in a form a program can read; the catalogue holds the window-level measurements from which the counts,
the coverage and the sensitivity of this paper are formed.

What a reader can and cannot repeat is worth saying plainly. The observations
can be obtained in full from the archive and the method is specified in enough
detail to be reimplemented, so the search can be carried out again
independently, which is the stronger test in any case. The calibrated
visibilities, the per-window products and the code that produced them are not
distributed with this paper, so this particular execution of the search cannot
be re-run by a reader. Two choices that could create or remove an unattributed
crossing were fixed before the data they govern were searched: the calibration
hold-out rule of Appendix~\ref{sec:heldout} and the block list its first guard
depends on were fixed \HoLeadBlockDays{} days before the first reserved block
was searched, and the disposition criteria were fixed on \MkCritDate, before
the line list was queried. The line list is a dated query of a public
catalogue and can be repeated exactly as it stands; the ordering of the other
two is a statement by the authors, and a reader has no way to check it.
